% id: 28-53
% book: 쎈B 미적분1 · 02 함수의 연속
% section: 유형 12 사잇값 정리의 실생활에의 활용
% figure: none
% answer: 
% answer_source: 
% variant_level: 0 (원본 전사)
% 이 파일은 build.mjs 가 items.json 에서 만든다. 고칠 때는 items.json 을 고친다.
\smash{\raisebox{1.5mm}{\dmkichul{쎈B 미적분1 28쪽 53번}}}
어느 날 $\pt{A}$ 지역의 11시부터 16시까지의 기온의 변화를 1시간 단위로 측정하였더니 아래 표와 같았다. \par\smallskip\begin{center}\small\begin{tabular}{|c|c|c|c|c|c|c|}\hline 시각 & 11시 & 12시 & 13시 & 14시 & 15시 & 16시 \\ \hline 기온($^{\circ}\mathrm{C}$) & $23$ & $24$ & $26$ & $26$ & $22$ & $23$ \\ \hline\end{tabular}\end{center}\smallskip\noindent 다음 중 이 지역의 기온에 대한 설명으로 항상 옳은 것은?
\par\medskip\par\noindent\hangindent=1.4em\hangafter=1 {\small ①}\ 11시부터 12시까지 기온이 $23.5\mkern3mu ^{\circ}\mathrm{C}$가 되는 때가 적어도 두 번 있었다.\par\noindent\hangindent=1.4em\hangafter=1 {\small ②}\ 12시부터 13시까지 기온이 $23\mkern3mu ^{\circ}\mathrm{C}$가 되는 때가 적어도 한 번 있었다.\par\noindent\hangindent=1.4em\hangafter=1 {\small ③}\ 13시부터 14시까지 기온이 $25\mkern3mu ^{\circ}\mathrm{C}$가 되는 때가 적어도 두 번 있었다.\par\noindent\hangindent=1.4em\hangafter=1 {\small ④}\ 13시부터 16시까지 기온이 $22.5\mkern3mu ^{\circ}\mathrm{C}$가 되는 때가 적어도 두 번 있었다.\par\noindent\hangindent=1.4em\hangafter=1 {\small ⑤}\ 11시부터 16시까지 기온이 $24.5\mkern3mu ^{\circ}\mathrm{C}$가 되는 때가 적어도 세 번 있었다.\par\medskip
